%% Currículo de exemplo - Marcos Santos
%% Visão completa de todas as competências, experiências e conquistas.
%% Use como referência mestre ao personalizar currículos por vaga.
%%
%% SETUP: rode o /setup para preencher com as suas informações reais.
%% Compilar com: lualatex main_example.tex
%% (o pdflatex costuma falhar no MiKTeX recente com erros de expansão de fonte do fontawesome5.)
%% Convenções brasileiras: sem foto, CPF, RG, idade ou estado civil; datas MM/AAAA - MM/AAAA;
%% valores em reais com escape: R\$ (veja 05-cv-templates.md).

\documentclass[11pt,a4paper,sans]{moderncv}
\moderncvstyle{banking}
\moderncvcolor{blue}

% Força o nome e os títulos de seção a saírem no azul do moderncv (color1).
% O banking padrão deixa os dois em preto: o \colorlet de
% moderncvstylebanking.sty copia a cor de destaque anterior ao esquema, então
% as cores do nome ficam congeladas antes do \moderncvcolor. Redefinir depois
% resolve. \namefont é o gancho de todas as macros de nome, então isto funciona
% também no moderncv 2.3.1 (apt do Debian/Ubuntu), sem \firstnamestyle/\lastnamestyle.
\renewcommand*{\namefont}{\fontsize{34}{36}\bfseries\upshape}
\colorlet{firstnamecolor}{color1}
\colorlet{lastnamecolor}{color1}
\colorlet{namecolor}{color1}
\renewcommand*{\sectionstyle}[1]{{\sectionfont\color{color1}#1}}

\usepackage[utf8]{inputenc}
% Só para o caso de compilar com pdflatex (o motor documentado é o lualatex,
% que pula este ramo). Sem a codificação T1, o pdflatex monta letras acentuadas
% com \accent e a camada de texto do PDF as guarda decompostas - `a` + U+0303 em
% vez de U+00E3 -, então um ATS que procura "Gestão" não encontra, embora a
% página pareça certa. O moderncv 2.5 carrega T1 sozinho no pdflatex; o 2.3.1
% (apt do Debian/Ubuntu) não. \ifpdftex vem do iftex, que todo moderncv carrega.
\ifpdftex\usepackage[T1]{fontenc}\fi
% O moderncv carrega o hyperref num gancho \AtEndPreamble, então o \hypersetup
% precisa ir num \AtEndPreamble nosso: no moderncv < 2.4 um \usepackage{hyperref}
% solto conflita com o \RequirePackage[unicode]{hyperref} da própria classe. A
% partir da 2.4.0 a classe repassa as opções por \PassOptionsToPackage, o que
% elimina o conflito.
\AtEndPreamble{\hypersetup{
    colorlinks=true,
    linkcolor=blue,
    filecolor=magenta,
    urlcolor=blue,
    pdftitle={Marcos Santos - CV},
    % Mantenha pdfpagemode=UseNone: este bloco roda depois do \AtEndPreamble
    % do próprio moderncv (que define pdfpagemode ali), então um FullScreen aqui
    % ganharia e abriria todo currículo em tela cheia.
    pdfpagemode=UseNone,
}}
\usepackage[scale=0.80]{geometry}
\usepackage{import}

% dados pessoais (veja as convenções do currículo brasileiro em 05-cv-templates.md)
\name{Marcos}{Santos}
% Basta "Cidade - UF". Sem endereço a listar, APAGUE a linha inteira:
% \address{}{}{} falha com "There's no line here to end" em todas as versões do moderncv.
\address{Lagoa da Prata - MG}{}{}
\phone[mobile]{+55 37 99829-2763}
\email{marcossantosjdev@gmail.com}
\extrainfo{\href{https://www.linkedin.com/in/marcos-santos-b9b544214/}{LinkedIn}, \href{https://github.com/Marcos-Vinicius-F-Santos}{GitHub}}

\begin{document}

\makecvtitle

% ============================================================
%     RESUMO PROFISSIONAL
% ============================================================

\vspace{6pt}
\small{Desenvolvedor de software com mais de 2 anos de experiência em aplicações corporativas, com foco em backend Java, integrações e entregas ponta a ponta. Experiência com Java 17/21, REST, SQL, SAP RFC/BAPI/JCo e Git, além de Spring Boot, testes automatizados, Flyway e Docker em projetos pessoais. Atua conectando levantamento com negócio, modelagem de dados, implementação e evolução de soluções com resultados mensuráveis.}

% ============================================================
%     COMPETÊNCIAS
% ============================================================

\section{Competências}
\vspace{1pt}
\begin{itemize}

\item \textbf{Backend e APIs}: Java 17/21, Spring Boot, Spring Data JPA, REST APIs, Flyway e modelagem de serviços.

\item \textbf{Dados e testes}: SQL, MySQL, PostgreSQL, JUnit, Mockito, Testcontainers e testes automatizados.

\item \textbf{Frontend}: React, Angular, TypeScript, JavaScript e HTML/CSS.

\item \textbf{Integrações e domínio}: SAP RFC/BAPI/JCo, aplicações corporativas, transformação digital, agronegócio e laticínios.

\item \textbf{Ferramentas}: Git, code review, Maven, Gradle, Docker, Docker Compose, Grafana, OutSystems, Mendix e AWS Fundamentals em andamento.

\end{itemize}

% ============================================================
%     EXPERIÊNCIA PROFISSIONAL
% ============================================================

\section{Experiência Profissional}
\vspace{3pt}
\begin{itemize}

% --- Cargo mais recente ---
\item{\cventry{05/2025 - 03/2026}{Junior Software Engineer}{Alvoar Lácteos}{Lagoa da Prata - MG}{}{\vspace{1pt}
\begin{itemize}
    \item {Conduzi uma aplicação corporativa crítica de onboarding ponta a ponta, do levantamento com negócio e modelagem de dados à implementação de frontend, backend e integrações.}
    \item {Reduzi etapas manuais de 8-10 para 3-5 e o onboarding de dias para horas, em um fluxo usado por mais de 2.000 usuários mensais.}
    \item {Desenvolvi integrações em Java 17/21 com REST APIs, SAP RFC/BAPI/JCo e bancos SQL; um componente reutilizável de desserialização JSON reduziu o esforço de mapeamento em aproximadamente 50\%.}
    \item {Defini estruturas de dados, APIs e módulos compartilhados com Git, code review, Scrum/Kanban, validações manuais e Grafana.}
\end{itemize}}}

\vspace{3pt}

% --- Cargo anterior ---
\item{\cventry{12/2023 - 05/2025}{Software Developer}{Alvoar Lácteos}{Lagoa da Prata - MG}{}{\vspace{1pt}
\begin{itemize}
    \item {Desenvolvi aplicações internas em OutSystems e componentes customizados com React e JavaScript, contribuindo para reduzir em 30\% o tempo de conclusão de tarefas.}
    \item {Implementei formulários dinâmicos que reduziram em 25\% os pedidos de modificação e trabalhei diretamente com usuários para refinar requisitos e validar entregas.}
\end{itemize}}}

\end{itemize}

% ============================================================
%     FORMAÇÃO ACADÊMICA
% ============================================================

\section{Formação Acadêmica}
\vspace{1pt}
\begin{itemize}

\item{\cventry{07/2023 - 05/2024}{Pós-graduação lato sensu em Business Intelligence, Big Data e Analytics - Ciência de Dados}{Unopar Anhanguera}{[Cidade - UF]}{}{\vspace{1pt}
[KEY_TOPICS]
}}

\vspace{3pt}

\item{\cventry{12/2023 - 06/2024}{Pós-graduação lato sensu em Estatística Aplicada}{Unopar Anhanguera}{[Cidade - UF]}{}{\vspace{1pt}
[KEY_TOPICS]
}}

\vspace{3pt}

\item{\cventry{01/2022 - 06/2023}{Tecnólogo em Design Gráfico}{Unopar}{[Cidade - UF]}{}{\vspace{1pt}
[KEY_TOPICS]
}}

\end{itemize}

% ============================================================
%     IDIOMAS
% ============================================================

\section{Idiomas}
\vspace{1pt}
\begin{itemize}
\item {Português (nativo), Inglês (C1/C2), Francês (A1).}
\end{itemize}

% ============================================================
%     PUBLICAÇÕES (opcional)
% ============================================================

\section{Publicações}
\vspace{1pt}
\begin{itemize}
\item {[Autoria] ([Ano]). [Título]. [Periódico/Congresso]. \href{[DOI_URL]}{link do DOI}}
\end{itemize}

% ============================================================
%     PRÊMIOS E RECONHECIMENTOS (opcional)
% ============================================================

\section{Prêmios e Reconhecimentos}
\vspace{1pt}
\begin{itemize}
\item{\textbf{[Nome do prêmio]} -- [Evento/Organização] ([Ano]).}
\end{itemize}

% ============================================================
%     REFERÊNCIAS (no Brasil, em geral omitida; mantenha só se sobrar espaço)
% ============================================================

\section{Referências}
\vspace{1pt}
\begin{itemize}
\item{Disponíveis mediante solicitação.}
\end{itemize}

\end{document}
